\scanpage{199}
\mspage{II.47}
% archive sheet number: 198
% lines 1--7
Dans chaque pb géométrique particulier, \unsure{on n'a guère qu'à} considérer les représentations (analytiques ou algébriques) se \unsure{factorisant|factorisent} par un gpe \deleted{quotient} $\unsure{G|C}_\alpha$ resp. $\widetilde G_\alpha$ fixé du syst. projectif, dont on peut espérer que c'est l'image dans un $\underline{\operatorname{Aut}}_{\mathbb{Q}_{\unsure{\ell|l}}}(E)$ d'un gpe $\pi_1(X_\lambda,\unsure{\xi|\eta}_\lambda)$, resp. son enveloppe algébrique.

% lines 8--14
Notons que \deleted{\unsure{pour tout le}} dans (6.15.1) l'opération de $\Pi$ sur $E$ définit un gpe algébrique $\widetilde G_E\subset\underline{\operatorname{Aut}}_{\mathbb{Q}_\ell}(E)$, enveloppe algébrique d'un gpe analytique $G_E$ \deleted{\unsure{provenant}} \inserted{image} d'un $\pi_1(X_\lambda,\xi)$, et que le deuxième \unsure{membre} de (6.15.1), peut aussi s'écrire, \unsure{même|mais} formellement, $\operatorname{End}_{G_E}(E)=\operatorname{End}_{\widetilde G_E}(E)$.

% lines 15--19
Si $\pi_1(X,\xi)$ (\uline{non} arith) $=1$ (p. ex $X$ le spectre d'un corps sep clos) alors ce deuxième membre est aussi $\operatorname{End}_{\unsure{\mathfrak g|g}_E}(E)=\operatorname{End}_{\widetilde{\mathfrak g}_E}(E)$, où $\mathfrak g_E=\operatorname{Lie}(G_E)\subset\operatorname{End}(E)$ et $\widetilde{\mathfrak g}_E$ est l'enveloppe algébrique de $\mathfrak g_E$ ....
